\chapter*{Abstract}
\addcontentsline{toc}{chapter}{Abstract}

%% Edit below this line
Sign language is the primary means of communication for millions of deaf and
hard-of-hearing people, yet the scarcity of interpreters in everyday settings
creates a persistent barrier with the hearing world. This project presents a
self-contained, wearable glove that recognises American Sign Language (ASL)
gestures and translates them into spoken output in real time, without a camera
or any external infrastructure.

The system places a compact inertial measurement unit (MPU6050) and a
microcontroller (Seeed Studio XIAO ESP32-C3) on the nail of each of the five
fingers. Four ``slave'' modules stream fused orientation and angular-rate data
over the ESP-NOW protocol to a fifth ``master'' module on the middle finger,
which aggregates a 45-channel feature vector at 25\,Hz and forwards it to a
mobile phone over Bluetooth Low Energy. A Flutter application captures a
fixed-length window of sensor data on demand, runs a quantised one-dimensional
convolutional neural network entirely on-device, and speaks the recognised
letter or word through text-to-speech.

A central finding of the work is the decisive role of the feature
representation. Static handshape recognition based on absolute finger
orientation proved fragile, because the inertial reference frame cannot be
reproduced exactly between sessions and because the limited angular resolution
of the sensors cannot separate handshapes that differ only by small thumb
displacements. The project therefore adopts distinctive per-letter
\emph{motions} together with \emph{absolute} (non-orientation-normalised)
features, relocating the discriminative information into the
calibration-independent angular-rate channels. The deployed models occupy
137\,KB of flash, require 22.7\,KB of RAM, and execute in approximately 1\,ms
per inference; the wireless link delivers all five finger streams at 25\,Hz
with negligible packet loss. The project demonstrates that accurate,
low-latency, fully on-device sign-language translation is achievable on
inexpensive commodity hardware, and it documents the engineering trade-offs—%
radio coexistence, power integrity, thermal behaviour, and feature
invariance—that govern such a system.

%% Until here
\vfill
%% Edit after {Keywords:}
\textbf{Keywords:} sign language recognition, wearable computing, inertial
measurement unit, ESP-NOW, Bluetooth Low Energy, edge machine learning,
one-dimensional convolutional neural network, gesture recognition.
\clearpage
